\section{直线题库}

\begin{problemset}
    \item 过点 $(-2,0)$ 与圆 $x^2+y^2-6 x=0$ 相切的两条直线的夹角为 $\alpha$ ，则 $\sin \alpha=$ \underline{$\qquad$} ．
    \item 两直线 $y=x-1$ 与 $x-2 y+1=0$ 的夹角为 \underline{$\qquad$} ．（结果用反三角表示）
    \item 已知直线过点 $P(2,2)$ 且倾斜角为 $135^{\circ}$ ，则点 $Q(-2,0)$ 到直线 $l$ 的距离为 \underline{$\qquad$} ．
    \item 若直线 $y=-a x-2$ 与连接 $P(-2,1), Q(3,2)$ 两点的线段有公共点，求实数 $a$ 的取值范围
    \item 下列说法不正确的有 （\qquad）
    \begin{itemize}
        \item 直线的倾斜角越大，斜率越大
        \item 若直线 $l$ 经过点 $A(1,0), ~ B(2, \tan \alpha)$ ．则直线 $l$ 的倾斜角是 $\alpha$
        \item 直线 $x \sin \alpha+y+2=0$ 的倾斜角 $\theta$ 的取值范围是 $\left[0, \frac{\pi}{4}\right] \cup\left[\frac{3 \pi}{4}, \pi\right)$
        \item 直线 $\frac{x}{2}-\frac{y}{3}=1$ 在 $y$ 轴上的截距是 3
    \end{itemize}
    \item 已知点 $A(0,1), B(4,3)$ ，线段 $A B$ 的垂直平分线在 $x$ 轴上的截距为 \underline{\qquad} ．
    \item 已知点 $P(2,-1)$ ，直线 $l: 3 x+2 y+5=0$ ．过点 $P$ 且与直线 $l$ 垂直，求直线 $l$ 的方程 \underline{\qquad}.
    \item 设 $a$ 为实数，直线 $l_1: a x+y=1$ ，直线 $l_2: x+a y=2 a$ ，则＂$a=1$＂是＂$l_1, l_2$ 平行＂的（\quad）条件
    \item "$m=-4$"是"直线 $l_1:(m-2) x-3 y-1=0$ 与直线 $l_2: m x+(m+2) y+1=0$ 相互平行"的 （\qquad）
\end{problemset}

\v{2em}

\subsection{练习答案}
\begin{tips}
    \begin{enumerate}
        \item 过点 $(-2,0)$ 与圆 $x^2+y^2-6 x=0$ 相切的两条直线的夹角为 $\alpha$ ，则 $\sin \alpha= \frac{24}{25}$ ．
        \item 两直线 $y=x-1$ 与 $x-2 y+1=0$ 的夹角为 $arctan\frac{1}{3}$
        \item 已知直线过点 $P(2,2)$ 且倾斜角为 $135^{\circ}$ ，则点 $Q(-2,0)$ 到直线 $l$ 的距离为 $3\sqrt{2}$ ．
        \item 已知直线 $l$：$(m+1) x+(2 m-1) y+m-2=0$ ，则直线恒过定点 \underline{$(1,-1)$}.
        \item 若直线 $y=-a x-2$ 与连接 $P(-2,1), Q(3,2)$ 两点的线段有公共点，实数 $a$ 的取值范围是 $\left(-\infty,-\frac{4}{3}\right] \cup\left[\frac{3}{2},+\infty\right)$
        \item 选择（ABD）
        \item $3$.
        \item $2 x-3 y-7=0$
        \item 充分不必要条件

              若 $a=1$，则直线 $l_1: x+y=1$，直线 $l_2: x+y=2$。此时两直线平行，所以充分性成立。

              若 $l_1 \parallel l_2$，则有 $a \cdot 1 = 1 \cdot 1$，解得 $a^2=1$，即 $a = \pm 1$。  当 $a=1$ 或 $a=-1$ 时，两直线均平行。

              因此，$l_1 \parallel l_2$ 不能唯一推出 $a=1$，必要性不成立。
              故为充分不必要条件。
        \item 充要条件

              两直线平行的条件是 $(m-2)(m+2) = -3m$，且截距不相等。

              化简得 $m^2+3m-4=0$，解得 $m=-4$ 或 $m=1$。

              当 $m=-4$ 时，两直线分别为 $-6x-3y-1=0$ 和 $-4x-2y+1=0$，斜率均为 $-2$，截距不同，互相平行。

              当 $m=1$ 时，两直线方程相同，为重合关系，不符合平行。

              因此，$m=-4$ 是两直线平行的充要条件。
    \end{enumerate}
\end{tips}

\v{12em}